\documentclass[a4paper]{article}
\usepackage[pdftex]{graphicx}
\usepackage[utf8]{inputenc}
\usepackage{enumerate}
\usepackage{amssymb}
\usepackage{icomma}
\usepackage{tikz}
\usepackage{href-ul}
\hypersetup{
	colorlinks=true,
	linkcolor=blue,
	urlcolor=blue}
\usepackage{geometry}
\geometry{a4paper, top=15mm, left=15mm, right=15mm, bottom=15mm,
	headsep=10mm, footskip=12mm}

% Start the document
\begin{document}
\begin{tikzpicture}(10,3)
	\draw[ultra thick](2,0) --(10,0) -- (10,1.3) --(2,1.3) -- (2,0);
	\draw[fill=black](2,0)-- (0,.6) -- (2,.6) -- (2,0);
	\node at (6,.6) {\large Lösungsblatt von \href{https://www.okuyakl.de}{www.okuyakl.de}};
\end{tikzpicture}
\vspace{0.5 cm}

\noindent{\bf Aufgabe 1.}\\
\renewcommand{\arraystretch}{3}
\begin{tabular}{|p{40 pt}|p{40 pt}|p{60 pt}|p{50 pt}|p{70 pt}|p{70 pt}|p{90 pt}|}
\hline
Funktion & Potenz & senkrechte Asymptote & waagrechte Asymptote  & Definitions\-menge & Wertemenge & Funktionsterm  \\
   & &  $x=$  &  $y=$ & $\mathbb{D}=$ & $\mathbb{W}=$ & $f(x)=$   \\
\hline
$f_1(x)$ &  $2$ & $-4$ & $-2$ & $\mathbb{D}=\mathbb{R} \setminus \{-4\}$ & $\mathbb{W}=]-2;\infty[$ & $f(x)={1\over (x+4)^2}-2$\\
\hline
$f_2(x)$ & $1$ & $-1$ & 0 & $\mathbb{D}=\mathbb{R} \setminus \{-1\}$ & $\mathbb{W}=\mathbb{R} \setminus \{0\}$ & $f(x)= {1 \over x+1}$\\
\hline
$f_3(x)$ & $1$ & $4$ & $-3$ & $\mathbb{D}=\mathbb{R} \setminus \{4\}$ & $\mathbb{W}=\mathbb{R} \setminus \{-3\}$ & $f(x)= {1 \over x+4}-3$\\
\hline
$f_4(x)$ & 2 & 5 & $0$ & $\mathbb{D}=\mathbb{R} \setminus \{5\}$ & $\mathbb{W}=]0;\infty[$ & $f(x)= {1 \over (x-5)^2}$\\
\hline
\end{tabular}

\begin{center}
	\includegraphics[width=7 cm]{../../viecher/endcomic.pdf}
	
	Hier geht es zurück zum \href{https://www.okuyakl.de/math/m8hypL019/aa019.pdf}{Aufgabenblatt}
\end{center}

\end{document}